% !PRINTED 392
\weberpage{392}
%<<CONT
$\omega_n$ eine quadratische Gleichung für $\omega$, die man in die Form
setzen kann:
\[
\begin{vmatrix} P_n - \omega Q_n, & P_{n-1} - \omega Q_{n-1} \\ P_m - \omega Q_m, & P_{m-1} - \omega Q_{m-1} \end{vmatrix} = 0.
\]

\begin{center}
\S.\ 126.

Beispiele.
\end{center}

Wir wollen einige Beispiele für die Bestimmung der reducirten
 Zahlen und Perioden hier durchführen, um die Anwendung
der Methode zu zeigen; die Beispiele lassen sich natürlich ganz
nach Belieben vermehren.

1. $D = 29$.

Wir bestimmen zunächst nach \S.\,125, 4. die sämmtlichen
zu $D$ gehörigen reducirten Zahlen. Die grösste in $\sqrt{D}$ enthaltene
ganze Zahl $\lambda$ ist hier $= 5$; also kann $b$ nur die Werthe haben:
\[
b = 1, 3, 5,
\]
woraus
\[
ac = \frac{D-b^2}{4} = 7, 5, 1.
\]

Die Grenzen, in denen $a$ und $c$ liegen müssen, sind nach
\S.\,125, (7) für die drei Werthe von $b$:
\[
3, 3; \quad 2, 4; \quad 1, 5.
\]
Man erhält also nur die einzige reducirte Zahl:
\[
\{1, 5, 1\}.
\]

Die Periode besteht aus einem einzigen Gliede, und der
Kettenbruch für $\omega$ ist:
\[
\omega = \frac{5+\sqrt{29}}{2} = (5, 5, 5 \ldots)
\]

2. $D = 116 = 4.29$.

Hier ist $\lambda = 10$ und $b$ hat einen der Werthe
\[
b = 2, 4, 6, 8, 10,
\]
also
\[
ac = 28, 25, 20, 13, 4;
\]
die Grenzen für $a$ und $c$ sind
\[
5, 6; \quad 4, 7; \quad 3, 8; \quad 2, 9; \quad 1, 10,
\]
also, da die Fälle, in denen $a$, $b$, $c$ den gemeinschaftlichen Factor
%CONT>>
